(1) 单单元配分函数 \(z=1+2e^{-\beta\Delta}+e^{-3\beta\Delta}=1+2q+q^3\)，因各单元独立，总配分函数 \(Z=z^N\)。单单元平均能量为 \(u=-\partial_\beta\ln z=\Delta(2q+3q^3)/z\)，故 \(U=Nu\)。两倍简并态的统计权重必须记为 \(2q\)，而非把三个不同能级一律赋相同权重。

(2) 利用正则系综涨落恒等式 \(C_V=N(\langle\epsilon^2\rangle-\langle\epsilon\rangle^2)/(k_BT^2)\)，得到 \(C_V=Nk_Bx^2\{(2q+9q^3)/z-[(2q+3q^3)/z]^2\}\)。方括号是能量以 \(\Delta\) 为单位的方差，因而 \(C_V\ge0\)，在 \(T>0\) 且能级不简并为同一能量时严格正。这样也避免直接对分式求温度导数可能产生的符号错误。

(3) 低温 \(x\gg1\) 时 \(z=1+2q+O(q^3)\)，\(u=2\Delta e^{-x}+O(e^{-2x})\)，所以 \(C_V=2Nk_Bx^2e^{-x}[1+o(1)]\)。高温 \(x\ll1\) 时四个微观态等概率，对 \(\epsilon/\Delta\) 的分布为 \((0,1,1,3)\)，平均 \(5/4\)，均方 \(11/4\)，方差 \(11/4-25/16=19/16\)，故 \(C_V\sim(19/16)Nk_Bx^2\)。高低温两端皆趋零，有限温度必出现热容峰。

(4) 熵由 \(S=k_B(\ln Z+\beta U)\) 给出，即 \(S=Nk_B[\ln z+x(2q+3q^3)/z]\)。低温 \(T\to0^+\) 时 \(S\to0\)，因为基态非简并；高温 \(q\to1\)、\(z\to4\)，每个局域单元有四个等概率微观态，故 \(S\to Nk_B\ln4\)。二重简并激发态是多出的一个微观态而不是新能级数目；不可把该极限写成 \(Nk_B\ln3\)。
